\documentclass[UTF8,11pt,a4paper]{ctexart}
\usepackage[margin=22mm]{geometry}
\usepackage{booktabs,tabularx,array,amsmath}
\usepackage[hidelinks]{hyperref}
\setlength{\parindent}{0pt}
\setlength{\parskip}{6pt}
\renewcommand{\arraystretch}{1.15}
\ctexset{section={format=\large\bfseries}}
\begin{document}
\begin{titlepage}\centering\vspace*{0.32\textheight}
{\Huge\bfseries Pi 无额外 Skill\\[12pt]基线测试报告}
\end{titlepage}
\section{结论摘要}
本报告对应统一测试集 v2，执行时间为 2026-09-27 至 2026-10-01（UTC）。
开发集 600 对、正式集 1200 对，共 \textbf{3600 条有效轨迹}，全部完成；另有 24 条 canary 准入，不混入成绩。
2026-10-01 完成交付前复审，1800 对初态和稳定后状态一致，冻结配置及权重哈希核验通过。

\begin{center}
\begin{tabular}{lrrrr}
\toprule
正式集范围 & 每方法条数 & 直接 VLA & Pi & 差值（百分点）\\
\midrule
全部 & 1200 & 444 (37.00\%) & 454 (37.83\%) & +0.83\\
标准四组 & 400 & 339 (84.75\%) & 338 (84.50\%) & -0.25\\
Pro 八组 & 800 & 105 (13.13\%) & 116 (14.50\%) & +1.38\\
\bottomrule
\end{tabular}
\end{center}
上表为\textbf{终态成功}。正式集整体配对差值的任务聚类 95\% 区间为
\textbf{[-1.42, +3.08] 个百分点}，包含零；不能据此宣称无额外 skill 的 Pi 带来稳定成功率提升，也不意味着两方法等效。

Pi 在 62 对中成功而直接 VLA 失败，52 对出现相反结果，净增 10 次成功。
Pi 平均动作从 330 降至 282.01 步（减少 14.54\%），但平均每条耗时从 13.50 秒增至 159.64 秒（11.82 倍）。
Pi 正式集误报完成 159/1200（13.25\%）；在 603 次声称完成中，159 次终态不成功（26.37\%）。
因此，动作更少不等于更快、更便宜或完成判断更可靠。

\section{固定方法与预算}
\textbf{B0：}同一冻结 VLA 接收完整任务指令，执行至动作预算。
\textbf{B1：}原生 Pi 0.85.1、\texttt{deepseek-flash} 高推理与原生视觉，调用同一 VLA 并自行决定结束。
VLA 为 \path{lerobot/pi05_libero_finetuned_v044}，不是 RPent 的 RLinf 权重。

Pi 保留 read/bash/edit/write，机器人工具仅为 robot\_observe、robot\_execute\_vla、robot\_finish。
skills 与 contextFiles 为空；无额外机器人 skill、恢复模块或成功真值反馈。双相机图像、本体状态和时间戳可观察。
每个动作块 32 步，每次最多 4 块；动作预算为 Spatial 220、Object 280、Goal 300、Long 520，Pro 沿用对应家族预算。
Pi 每条墙钟上限 30 分钟，单次输出上限 65,536 token。环境/策略种子为 7。

真值只用于私有评分，不触发提前停止。B0 可能先成功后破坏目标，因此这些结果不能直接等同于首次成功即停止的论文成绩。
时间为记录中的每条执行墙钟耗时，含 API 等待与原生重试；不是 GPU 纯推理时间，也不是整个任务历时。
\newpage
\section{正式集：逐组结果}
每组每方法 100 条（10 个任务、每任务 10 个初态）。曾达成为任一物理步成功；终态为结束时成功。
误报仅针对 Pi 的完成声明，B0 无此声明。费用为美元估算，详见费用页。
\begin{center}\small
\begin{tabular}{llrrrrrr}
\toprule
环境 & 方法 & 曾达成 & 终态 & 误报 & 平均步 & 平均秒 & API估算\\
\midrule
Spatial & B0 & 94 & 80 & -- & 220.0 & 11.6 & 0\\
 & B1 & 92 & 88 & 10 & 150.8 & 97.3 & 9.69\\
Object & B0 & 98 & 98 & -- & 280.0 & 11.6 & 0\\
 & B1 & 97 & 97 & 0 & 187.7 & 110.6 & 14.91\\
Goal & B0 & 94 & 83 & -- & 300.0 & 11.4 & 0\\
 & B1 & 89 & 84 & 5 & 190.6 & 106.3 & 12.49\\
Long & B0 & 96 & 78 & -- & 520.0 & 18.9 & 0\\
 & B1 & 90 & 69 & 19 & 390.6 & 195.3 & 27.72\\
\midrule
Spatial Task & B0 & 36 & 35 & -- & 220.0 & 11.8 & 0\\
 & B1 & 37 & 37 & 10 & 194.8 & 159.4 & 20.69\\
Spatial Swap & B0 & 27 & 19 & -- & 220.0 & 12.0 & 0\\
 & B1 & 26 & 26 & 25 & 195.6 & 141.9 & 18.91\\
Object Task & B0 & 10 & 10 & -- & 280.0 & 12.2 & 0\\
 & B1 & 10 & 10 & 10 & 268.2 & 205.5 & 31.48\\
Object Swap & B0 & 5 & 5 & -- & 280.0 & 12.3 & 0\\
 & B1 & 9 & 9 & 11 & 273.0 & 194.8 & 25.72\\
Goal Task & B0 & 17 & 5 & -- & 300.0 & 11.7 & 0\\
 & B1 & 13 & 2 & 21 & 289.2 & 155.6 & 21.15\\
Goal Swap & B0 & 19 & 17 & -- & 300.0 & 11.5 & 0\\
 & B1 & 19 & 19 & 14 & 265.1 & 127.4 & 15.03\\
Long Task & B0 & 18 & 14 & -- & 520.0 & 18.6 & 0\\
 & B1 & 15 & 12 & 29 & 469.5 & 219.5 & 34.43\\
Long Swap & B0 & 0 & 0 & -- & 520.0 & 18.5 & 0\\
 & B1 & 4 & 1 & 5 & 508.8 & 201.9 & 28.05\\
\bottomrule
\end{tabular}
\end{center}
整体曾达成：B0 514/1200（42.83\%），B1 501/1200（41.75\%）。
整体终态：B0 444/1200，B1 454/1200。两者指标方向不同，不能仅选取有利指标。
Pi 的 Spatial 标准组终态较高、Long 标准组较低；这些分组差异为描述性结果，未进行多重比较显著性声明。
\newpage
\section{开发集：逐组结果}
每组每方法 50 条（10 个任务、每任务 5 个初态）。开发成绩独立报告，不与正式集混合。
\begin{center}\small
\begin{tabular}{llrrrrrr}
\toprule
环境 & 方法 & 曾达成 & 终态 & 误报 & 平均步 & 平均秒 & API估算\\
\midrule
Spatial & B0 & 46 & 42 & -- & 220.0 & 12.1 & 0\\
 & B1 & 45 & 45 & 3 & 152.6 & 82.2 & 4.44\\
Object & B0 & 49 & 49 & -- & 280.0 & 12.1 & 0\\
 & B1 & 50 & 49 & 1 & 186.6 & 86.0 & 5.92\\
Goal & B0 & 49 & 42 & -- & 300.0 & 11.6 & 0\\
 & B1 & 39 & 35 & 6 & 203.5 & 86.3 & 6.02\\
Long & B0 & 47 & 40 & -- & 520.0 & 20.0 & 0\\
 & B1 & 44 & 36 & 3 & 383.0 & 172.2 & 13.77\\
\midrule
Spatial Task & B0 & 16 & 15 & -- & 220.0 & 12.0 & 0\\
 & B1 & 15 & 15 & 12 & 191.1 & 146.6 & 11.04\\
Spatial Swap & B0 & 12 & 10 & -- & 220.0 & 12.3 & 0\\
 & B1 & 10 & 8 & 16 & 197.2 & 129.4 & 9.80\\
Object Task & B0 & 5 & 5 & -- & 280.0 & 12.2 & 0\\
 & B1 & 5 & 5 & 7 & 265.0 & 185.9 & 16.78\\
Object Swap & B0 & 3 & 3 & -- & 280.0 & 12.2 & 0\\
 & B1 & 7 & 7 & 10 & 271.4 & 190.0 & 15.23\\
Goal Task & B0 & 7 & 1 & -- & 300.0 & 11.8 & 0\\
 & B1 & 9 & 1 & 13 & 288.1 & 146.9 & 11.87\\
Goal Swap & B0 & 11 & 11 & -- & 300.0 & 11.5 & 0\\
 & B1 & 6 & 6 & 4 & 274.0 & 122.4 & 9.95\\
Long Task & B0 & 10 & 7 & -- & 520.0 & 19.0 & 0\\
 & B1 & 12 & 12 & 11 & 461.1 & 178.3 & 14.40\\
Long Swap & B0 & 1 & 0 & -- & 520.0 & 19.0 & 0\\
 & B1 & 1 & 0 & 4 & 507.5 & 177.3 & 13.27\\
\bottomrule
\end{tabular}
\end{center}
整体终态：B0 225/600（37.50\%），B1 219/600（36.50\%）。
曾达成：B0 256/600（42.67\%），B1 243/600（40.50\%）。
Pi 误报 90/600（15.00\%），声称完成 304 次中的误报占 29.61\%。
平均动作 B0 330、B1 281.76；平均秒数 B0 13.81、B1 141.96。
\newpage
\section{配对统计与不确定性}
将同一任务、扰动组与初态的两方法配对。这里的“救回”只表示 B0 失败而 B1 成功的配对差异；“损害”表示相反结果，\textbf{不是在同一条失败轨迹上实施干预后的因果恢复率}。

\begin{center}\small
\begin{tabular}{llrrrrl}
\toprule
划分 & 范围 & 对数 & 救回 & 损害 & 差值pp & 95\% 聚类区间pp\\
\midrule
开发 & 全部 & 600 & 33 & 39 & -1.00 & [-5.17, 2.50]\\
 & 标准 & 200 & 14 & 22 & -4.00 & [-11.50, 2.50]\\
 & Pro & 400 & 19 & 17 & +0.50 & [-2.75, 4.00]\\
正式 & 全部 & 1200 & 62 & 52 & +0.83 & [-1.42, 3.08]\\
 & 标准 & 400 & 29 & 30 & -0.25 & [-4.50, 4.00]\\
 & Pro & 800 & 33 & 22 & +1.38 & [-0.75, 3.50]\\
\bottomrule
\end{tabular}
\end{center}
按 40 个基础任务（家族、任务编号）聚类，Task/Swap 变体与所有初态保留在同一簇。
对簇有放回抽样 20,000 次，种子 20261001，取配对终态差值的 2.5\% 与 97.5\% 分位数。
各簇样本数平衡，因此簇均值与逐对均值一致。区间反映本测试集任务间差异，\textbf{不覆盖 API 多次随机重复、其他任务分布或真机迁移的不确定性}。
正式集两者均成功 392 对、均失败 694 对；开发集分别为 186 与 342 对。

\section{测试集与解释边界}
12 组由四个标准家族及各自 Task/Swap 扰动组成。总体分数按本清单加权：标准占 1/3、Pro 占 2/3，不代表任意实际部署任务分布。
开发初态为 0、1、2、47、48；正式初态为 7、44、17、11、24、27、14、38、49、32。
正式与开发初态不重叠，\textbf{但不是 VLA 未见任务测试}，也不是 1200 个独立任务。

本轮没有按成功率调整提示、工具、模型或动作预算。标准输入采用官方 metadata 指令，Pro 输入采用实际扰动 BDDL。
v1 曾在指令来源校验阶段失败，未执行模型；v2 仅修正标准指令来源，保留 v1 记录。
该结果只覆盖本次最小工具 Pi；不代表加入 skill、额外感知或成熟机器人 Agent 后的上限。

为后续迭代建立的分层清单为每 Agent 24 条接口检查、240 条日常开发、600 条候选正式评估、1200 条最终正式评估。
本次完整基线未缩减。候选 600 条属于最终 1200 条的子集，不能当成独立重复验证；正式子集不用于反复调参。复用本轮 B0 需要相同样本、模型、环境与预算。
\newpage
\section{费用与基础设施记录}
\begin{center}
\begin{tabular}{lrrrr}
\toprule
用途 & API 请求 & 正常返回 & 不确定/错误 & 美元估算含预留\\
\midrule
开发有效 Pi & 17559 & 17523 & 36 & 132.4927\\
正式有效 Pi & 34707 & 34642 & 65 & 260.2798\\
Canary 准入 & 450 & 450 & 0 & 3.3881\\
归档无效原轨迹 & 96 & 77 & 19 & 7.7993\\
健康检查 & 2 & 2 & 0 & 0.000075\\
\midrule
合计 & 52814 & 52694 & 120 & 403.9599\\
\bottomrule
\end{tabular}
\end{center}
按归档的输入 0.30、输出 1.20 美元/百万 token 计算，全部输入按缓存未命中、峰时处理。
2026-10-01 查阅 DeepSeek 官方定价页，确认该峰时价格仍与归档一致；官方另有更低的缓存命中及非峰时价格。
\textbf{以上为保守估算，不是账单实付，也不是严格保证的实际账单上界。}
正常返回请求估算 358.5227 美元，120 次不确定/错误请求保留 45.4372 美元预留；无未决 reserved 项。
不确定请求的实际扣费未核实，不静默按零处理。没有获取账户账单；不含 GPU、服务器与电费。
有效轨迹费用包含恢复后的替代轨迹，旧无效轨迹另列，避免遗漏或重复归入成绩。

官方价格来源（核查日期 2026-10-01）：
\url{https://api-docs.deepseek.com/quick_start/pricing/}

\textbf{接口故障与恢复：}
开发 Long task3/init0 的 Pi 原轨迹，以及正式 Object Task task7/init17 的 Pi 原轨迹，因上游连接错误未有效结束，分别独立版本化重测。
重测沿用原模型、提示、预算与初态；核对有效性和初态哈希后替换规范结果，完整旧轨迹及费用保留在 invalid-attempts。
仅有这两条作为无效原轨迹归档，未把未运行或接口无效条目算作 Agent 能力失败。

开发结束后，原进程因第二次环境 bootstrap 保护性退出，正式集尚未执行；改用新进程调用原冻结流程继续。
没有修改策略或重封配置。原生自行恢复的网络错误仍保留在有效轨迹的耗时与费用里，因此本报告时延包含基础设施影响。

Pi 同轮并行提交有依赖的裁剪与读取所致 ENOENT，以及未承诺依赖的探测失败，没有据此改策略或补装能力。
运行有效不等于每一个模型动作决策正确；低成功率也不单独归因于模型或工具的某一环节。
\newpage
\section{审计与可追溯归档}
复审逐条检查清单覆盖与唯一性、每步私有成功真值和目标谓词一致性、动作步数与预算、配对初态及稳定状态哈希。
1800 条有效 Pi 轨迹均记录图像输入（开发 600、正式 1200）；检查原生视觉请求与工具推理历史按原始返回内容回传。
检查 skills/contextFiles 为空、工具集合恰为四个原生工具及三个机器人工具、高推理配置不变。
Canary 已人工核对 12 组初末双相机与物理推进；全量审计不等于人工逐帧语义标注所有视频。

API 请求与返回名称均为 \texttt{deepseek-flash}；服务商可变别名无法冻结实际权重。
文件系统 jail 不是完整网络隔离。审计支持既定接口与记录完整性，不构成对所有可能信息通道的形式化安全证明。

主实验位置：
{\small\path{/root/autodl-tmp/experiments/pi-baseline-20260927-v2}}

development/ 与 evaluation/ 保存逐条结果、公共执行记录及私有评分记录；api/ 保存请求、响应和账本；invalid-attempts/ 保存无效原轨迹及替代映射。
统计脚本和统计 JSON 与报告一起归档，可从 audit JSON 重算配对区间。

\begin{tabularx}{\linewidth}{lX}
\toprule
文件 & SHA256\\
\midrule
protocol-lock.json & \small\path{4640b3f0e881e03d2614c2e154b2a955062139253195f3727e40512773464f5b}\\
development-audit.json & \small\path{c4317da7c90b813d2f559148dcd71053a6780cba601e7d558ca8f2b1fc1bfcfb}\\
evaluation-audit.json & \small\path{c0027443359bc06f13fc4a3e8b4a36f100f6ee8cbfe1ee50f23d60bb3befccef}\\
ledger.json & \small\path{eee34028d9fd9db44931fc27003d5f525b10b7f45ab2415372a13b715ed43a69}\\
\bottomrule
\end{tabularx}

本报告替换旧 140 条基线报告，旧版先归档；不混入成熟 Agent 对照报告。
仅同步本地与指定服务器，不推送 GitHub，不删除研究数据，不关闭服务器。
\end{document}
